\section{Discussion and Concluding Remarks}
\label{sec:discussion}

This work provides a multidimensional characterization of the effects of the weight decay hyperparameter within the modern language-model training lifecycle. While traditional perspectives have primarily viewed weight decay through the lenses of capacity control in over-parameterized regimes or optimization stability in single-epoch pretraining~\citep{zhang2017understanding, andriushchenko2023whywd}, our findings suggest that weight decay plays a far more nuanced role in shaping model behavior. In particular, we showed that models with smaller weight decay achieve lower validation loss after pretraining (especially in the over-trained regime), but that models with larger weight decay benefit from improved plasticity, enabling them to perform best when fine-tuned on downstream tasks. Weight decay may shape model plasticity through several mechanisms, including promoting linearly separable representations, regularizing attention matrices, and reducing overfitting on the training data.
Together, these findings reveal fundamental trade-offs in hyperparameter optimization. They also provide one of the first rigorous empirical demonstrations that selecting pretraining hyperparameters based solely on minimal pretraining validation loss can fail to yield the model with the highest performance on downstream tasks. 

The trade-offs we show mean that, in practice, the benefits of increased plasticity must be weighed against other effects that may depend on model size, training duration, and other parameters of the training setup. In heavily overtrained scenarios or for very large models trained for many steps \citep{singh2025openai,comanici2025gemini,anthropic2025system}, the benefits of markedly lower pretraining validation loss may outweigh those of plasticity. 
In addition, weight decay's diverse roles in training dynamics -- from plasticity (shown in this work) to optimization, training stability, convergence rate, and overfitting \citep{hoffmann2022training,andriushchenko2023whywd,kosson2025weight} -- adds further complexity to model training decisions. Its optimal value for one objective can conflict with that for another, as observed when optimizing for pretraining versus downstream performance. A single weight decay value may not satisfy multiple objectives, requiring weighing trade-offs and prioritizing objectives.


Future work may investigate in more detail the trade-offs between stability and plasticity, and the extent to which our results hold in large-model and heavily overtrained scenarios. They may also investigate the role of weight decay in model plasticity for foundation models beyond language (e.g., vision and multimodal foundation models) and for other downstream desiderata (beyond CoT reasoning, language understanding and commonsense reasoning, and safety alignment).
Taken together, the findings in this work cast light on the importance of model plasticity, the complexity of hyperparameter tuning throughout the training process of modern language models, and the multifaceted role that a single optimization hyperparameter plays in shaping model behavior.

